% Copyright (c) 2022 by Lars Spreng
% This work is licensed under the Creative Commons Attribution 4.0 International License. 
% To view a copy of this license, visit http://creativecommons.org/licenses/by/4.0/ or send a letter to Creative Commons, PO Box 1866, Mountain View, CA 94042, USA.

%~~~~~~~~~~~~~~~~~~~~~~~~~~~~~~~~~~~~~~~~~~~~~~~~~~~~~~~~~~~~~~~~~~~~~~~~~~~~~~
% You can add your packages and commands to the loadslides.tex file. 
% The files in the folder "styles" can be modified to change the layout and design of your slides.
% I have included examples on how to use the template below. 
% Some of these examples are taken from the Metropolis template.
%~~~~~~~~~~~~~~~~~~~~~~~~~~~~~~~~~~~~~~~~~~~~~~~~~~~~~~~~~~~~~~~~~~~~~~~~~~~~~~


\documentclass[
11pt,notheorems,hyperref={pdfauthor=whatever}
]{beamer}

\input{loadslides.tex} % Loads packages and some defined commands

\title[
% Text entered here will appear in the bottom middle
]{Ensemble learning}

\subtitle{How to use multiple times the same data}

\author[
% Text entered here will appear in the bottom left corner
]{
    Michele Vincenzo Branca, Leonardo De Pietro
}

\institute{
    Collegio Carlo Alberto, \\
    University of Turin}
\date{\today}

% packages
\usepackage{amsmath}

\begin{document}

% Generate title page
{
\setbeamertemplate{footline}{} 
\begin{frame}
  \titlepage
\end{frame}
}
\addtocounter{framenumber}{-1}

% You can declare different parts as a parentof sections
\begin{frame}{Index of the presentation}
    \tableofcontents[part=1]
    \tableofcontents[part=2]
    \tableofcontents[part=3]
\end{frame}

\makepart{What is Ensemble Learning?}

\justifying
\section{what is Ensemble Learning?}
\subsection{An overview}
\begin{frame}
    Given a \emphasize{\textbf{supervised learning algorithm}} \(\mathcal{A}(\mathcal{D})\,:\,\mathcal{X}\to\mathcal{Y}\), can we run it more than once on datasets constructed from the original one \(\mathcal{D}\) and then combine the results to get an overall better predictor? Kind of.
    
    First of all, what kind of combinations are we (typically) talking about?
    \begin{itemize}
        \item linear ones!
    \end{itemize}
    How do we do it for various algorithms?
    \begin{itemize}
        \item \alert{\textbf{local averaging methods}}: combine the predicted labels from estimators learned on different datasets;
        \item \alert{\textbf{regression}}: linearly combine the predictions;
        \item \alert{\textbf{classification}}: weighted majority vote or linear combination of real-valued predictions when convex surrogates are used.
    \end{itemize}
    What's the impact?
    \begin{itemize}
    \item for \example{\textbf{linear models}} this method does not add new functions;
    \item for \example{\textbf{non-linear models}} this method can add new functions. 
\end{itemize}
\end{frame}

\subsection{How to construct the datasets?}
\begin{frame}
    Suppose that \(\mathcal{D}\) is the original dataset. 

    We construct new datasets by assigning different weights to the data points. According to the weights' nature, they have different interpretations:
    \begin{itemize}
        \item \alert{\textbf{integer-valued weights}}: the new dataset is obtained by including the points the number of times specified by the weights;
        \item \alert{\textbf{(non-negative) real-valued weights}}: implemented by sampling points proportionally to their weights during training, effectively treating the normalized weight as a sampling probability (e.g., in \example{\textbf{SGD}});
        \item \alert{\textbf{arbitrary weights}}: modern techniques based on empirical risk minimization directly accomodate arbitrary weights. 
    \end{itemize}
    We will focus on \example{\textbf{bagging/averaging}} techniques:
    \begin{itemize}
        \item datasets constructed in parallel with \alert{\textbf{random weights}} (typically uniform) or - which yields similar results - using \alert{\textbf{random projections}}.
    \end{itemize}
\end{frame}

\subsection{A spoiler: what are the computational and statistical benefits?}
\begin{frame}
    It really depends on the original predictors:
    \begin{itemize}
        \item \alert{\textbf{local averaging methods}}: ensemble learning will typically \example{\textbf{reduce the variance}} of the predictor and parallalelization \example{\textbf{reduces the computational complexity}};
        \item \alert{\textbf{empirical risk minimization with nonlinear models}}: linear combinations \example{\textbf{increase the set of models}} from \(\varphi(\cdot,\,w)\) to \(\int\limits_{\mathcal{W}}\varphi(x,\,w)\,\mathrm{d}\nu(w)\), where \(\nu\) is a signed measure on \(\mathcal{W}\). Furthermore, \example{\textbf{variance and computational cost are decreased}};
        \item \alert{\textbf{empirical risk minimization with linear models}}: a linear combination of linear models will still be a linear models, thus not increasing the set of models. However, there are usually great \example{\textbf{computational benefits}}.
    \end{itemize} 
\end{frame}

\makepart{Averaging/Bagging}
\section{Averaging/Bagging}
\subsection{The framework}
\begin{frame}
    We will consider \alert{\textbf{regression with square loss}} where we have an explicit \example{\textbf{variance/bias}} decomposition.
    \begin{itemize}
        \item most results extend to other situations using convex surrogates.
    \end{itemize}
\end{frame}

\subsection{Averaging: the idealized setting}
\begin{frame}
    Suppose we have \(m\) indipendent datasets of size \(n\), all coming from an i.i.d. sample having the same distribution \(p(x,\,y)\) on \(\mathcal{X}\times\mathcal{Y}\). From each of them, we get an estimator:
    \[
        \hat{f}^{(j)},\quad j=1,\dots,\,m,
    \]
    and the new predictor is:
    \[
        \hat{f}^{\text{bag}}(x) = \frac{1}{m} \sum_{j=1}^m \hat{f}^{(j)}(x).
    \]
    Denote:
    \begin{itemize}
        \item \(\text{bias}^{(j)}(x)=\mathbb{E}\left[\hat{f}^{(j)}(x)\right]-f_*(x)\);
        \item \(\text{var}^{(j)}(x)=\text{var}\left[\hat{f}^{(j)}(x)\right]\);
    \end{itemize}
    where we assume that \(x\) is fixed, thus the expectations are taken with respect to the data only. Note that:
    \begin{itemize}
        \item \(\text{bias}^{\text{bag}}(x)=\text{bias}^{(j)}(x)\) for any \(i=1,\dots,\,m\);
        \item \(\text{var}^{\text{bag}}(x)=\text{bias}^{(j)}(x)\) for any \(i=1,\dots,\,m\);
    \end{itemize}
    because of i.i.d. assumption. 
\end{frame}

\subsection{Example: \(k\)-nearest neighbor regression}
\begin{frame}
    Take \(\mathcal{X}\subset\mathbb{R}^d\). In this case we have that the expected excess risk is upper bounded in this way\footnote{See Proposition 6.2, p. 170 \textit{Learning Theory from First Principles}, by \textit{Francis Bach}.}:
    \[
        \int_{\mathcal{X}}\mathbb{E}\left[(\hat{f}(x)-f_\star(x))^2\right]\mathrm{d}p(x)\leq\underbrace{\frac{\sigma^2}{k}}_{\text{variance upper bound}}+\underbrace{8B^2\text{diam}(\mathcal{X})^2\left(\frac{2k}{n}\right)^{2/d}}_{\text{bias upper bound}},
    \]
    where:
    \begin{itemize}
        \item \(\mathcal{X}\) is the \alert{\textbf{input feature space}} and \(\mathcal{Y}\) is the \alert{\textbf{output space}};
        \item \(f_\star\in\argmin\limits_{y\in\mathcal{Y}}\mathbb{E}\left[\ell(y,\,z)\mid x\right]\) is the \alert{\textbf{Bayes predictor}};
        \item \(\hat{f}(x)\) is the \alert{\textbf{prediction}} made by the \alert{\textbf{\(k\)-nearest neighbor estimator}} at a specific \alert{\textbf{evaluation point}} \(x\);
        \item \(p(x)\) is the \alert{\textbf{marginal probability distribution}} of the input features;
        \item \(\sigma^2\) is the \alert{\textbf{irreducible error}} \(\text{var}(y \mid x)\);
        \item \(B\) is the \alert{\textbf{Lipschitz constant}} of the Bayes predictor \(f_\star(x)\);
        \item \(\text{diam}(\mathcal{X})\) is the \alert{\textbf{maximum possible distance}} between any two points in the space.
    \end{itemize}   
\end{frame}
\begin{frame}    
    Thus, \alert{\textbf{averaging}} over \(m\) replications, we get that the \example{\textbf{excess risk}} is upper bounded by:
    \[
        \int_{\mathcal{X}}\mathbb{E}\left[(\hat{f}(x)-f_\star(x))^2\right]\mathrm{d}p(x)\leq\frac{\sigma^2}{\textcolor{red}{m}k}+8B^2\text{diam}(\mathcal{X})^2\left(\frac{2k}{n}\right)^{2/d},
    \]
    and optimizing over \(k\) yields: \(k\propto \frac{1}{m^{d/(2+d)}}\cdot n^{2/(2+d)}\). The overall excess risk ends up being proportional to \(1/(mn)^{2/(2+d)}\), which coincides with the rate for a dataset with \(N=nm\) observations.
    
    Clearly we have a limit concerning the number  of datasets we can construct: we need \(k>1\), thus substituting in this condition the optimal \(k\) (which is a function of \(m\) and \(n\)) we see that \(m\) cannnot grow larger than \(n^{2/d}\).
\end{frame}

\subsection{Example: Ridge regression (misspecified model)}
\begin{frame}
    In this case we have that the expected excess risk is upper bounded in this way\footnote{See Proposition 7.8, p. 216 \textit{Learning Theory from First Principles}, by \textit{Francis Bach}.}:
    \[
        \int_{\mathcal{X}}\mathbb{E}\left[(\hat{f}(x)-f_\star(x))^2\right]\mathrm{d}p(x)\leq\underbrace{\frac{\sigma^2}{n}\left(1+\frac{R^2}{\lambda n}\right)\text{tr}\left(\left(\Sigma+\lambda I\right)^{-1}\Sigma\right)}_{\text{variance upper bound}}+\underbrace{\left(1+\frac{R^2}{\lambda n}\right)^2\inf\limits_{f\in\mathcal{H}}\{\mid\mid f-f_\star\mid\mid^2_{L_2(p)}+\lambda\mid\mid f\mid\mid_{\mathcal{H}}^2\}}_{\text{bias upper bound}},
    \]
    where:
    \begin{itemize}
        \item \(\mathcal{X}\), \(f_\star\), \(\hat{f}\), \(p(x)\), \(\sigma^2\) represent the same things as before,
        \item \(\lambda\) is the \alert{\textbf{regularization parameter}}, which controls the trade-off between fitting the training data and keeping the model weights small;
        \item \(\Sigma\) is the \alert{\textbf{covariance matrix}} of the features in the feature space;
        \item \(R^2\) is a \alert{\textbf{bounding constant}} related to the norm of the features;
        \item \(\mathcal{H}\) is the \alert{\textbf{Reproducing Kernel Hilbert Space (RKHS)}} (or hypothesis space) where the candidate functions \(f\) live;
        \item \(\text{tr}\left(\left(\Sigma+\lambda I\right)^{-1}\Sigma\right)\) are the \alert{\textbf{degrees of freedom}} associated to the covariance operator \(\Sigma\).
    \end{itemize}  
\end{frame}

\begin{frame}    
    For the \example{\textbf{bias term}}, assuming:
    \begin{itemize}
        \item \(\mathcal{X}=\mathbb{R}^d\);
        \item that the target function belongs to a Sobolev Kernel of order \(t>0\);
        \item that \(\mathcal{H}\) is a Sobolev space of order \(s>d/2\); 
    \end{itemize}
    one gets that the bias term is of order \(\left(1+\frac{R^2}{\lambda n}\right)^2\lambda^{t/s}\). 
    
    For the \example{\textbf{variance term}}, assuming:
    \begin{itemize}
        \item that the eigenvalues \((\lambda_m)\) of the covariance matrix \(\Sigma\) satisfy \(\lambda_m\leq C(m+1)^{\alpha}\), \(\alpha>1\);
    \end{itemize} 
    then \(\text{tr}\left(\left(\Sigma+\lambda I\right)^{-1}\Sigma\right)\leq O(\lambda^{1/\alpha})\). With our \(\mathcal{X}\), \(\alpha=2s/d\), thus the variance is proportional to \(\frac{1}{n}\lambda^{-d/(2s)}\)\footnote{See section 7.6.6, p. 217 \textit{Learning Theory from First Principles}, by \textit{Francis Bach}.}.

    \alert{\textbf{Averaging}} over \(m\) replications, we get that the \example{\textbf{excess risk}} is proportional to \(\frac{\sigma^2}{nm}\lambda{-1/\alpha}+\lambda^{t/s}\), and again the estimator behaves like having \(N=nm\) observations.
\end{frame}

\makepart{The Two Paradigms: Sketching vs. Projection}

\justifying
\section{The Two Paradigms}
\subsection{Overview of Data Compression}
\begin{frame}
    In ordinary least-squares, when dealing with massive datasets, exact solutions require \(\mathcal{O}(nd^2)\) time. How can we use \emphasize{\textbf{random projections}} to speed this up or regularize the model?
    
    We consider two fundamental settings:
    \begin{itemize}
        \item \alert{\textbf{Sketching (\(S \in \mathbb{R}^{s \times n}\))}}: shrinking the datapoints (rows). We compress \(n\) observations into \(s\) new observations (\(n > s > d\)) by replacing \(\min_{\theta \in \mathbb{R}^d}\|y - \Phi\theta \|^2_2 \) by \(\min_{\theta \in \mathbb{R}^s}\|Sy - S\Phi\theta \|^2_2 \) where \(S\) is an i.i.d Gaussian matrix.
        \item \alert{\textbf{Random Projection (\(S \in \mathbb{R}^{d \times s}\))}}: shrinking the features (columns). We constrain our search space to an \(s\)-dimensional subspace (\(d > n > s\)) by replacing \(\min_{\theta \in \mathbb{R}^d}\|y - \Phi\theta \|^2_2 \) by \(\min_{\eta \in \mathbb{R}^s}\|y - \Phi S \eta \|^2_2 \).
    \end{itemize}
    What are the primary benefits?
    \begin{itemize}
    \item \example{\textbf{Sketching}} provides computational and storage benefits (reducing \(\mathbb{R}^{n \times d}\) to \(\mathbb{R}^{s \times d}\)).
    \item \example{\textbf{Random Projections}} provide both a reduction in computation time and inherent \example{\textbf{regularization}} in high-dimensional setups. 
\end{itemize}
\end{frame}

\subsection{Connection to Random Forests}
\begin{frame}
    How do these concepts map to well-known algorithms?
    
    \emphasize{\textbf{Random Forests}} represent the ultimate unification of these two dimensionality reduction concepts in classical machine learning:
    \begin{itemize}
        \item \alert{\textbf{Bagging/Subsampling}}: Every decision tree is trained on a bootstrapped sample of rows (analogous to Sketching).
        \item \alert{\textbf{Feature Selection}}: At every node split, the algorithm selects a random subset of features (analogous to Random Projections).
    \end{itemize}
    This dual-randomization maximally decorrelates individual estimators, driving the \example{\textbf{variance term}} of the ensemble down far more effectively than bagging alone.
\end{frame}

\makepart{Gaussian Sketching}
\section{Gaussian Sketching}
\subsection{The Framework}
\begin{frame}
    Following the ordinary least-squares framework, we consider a design matrix \(\Phi \in \mathbb{R}^{n \times d}\) with full column rank \(d\) (thus \(n \ge d\)). 
    
    We generate \(m\) replications using Gaussian sketching matrices \(S^{(j)} \in \mathbb{R}^{s \times n}\). The \alert{\textbf{sketched closed-form estimator}} is:
    \[
        \hat{\theta}^{(j)} = (\Phi^\top(S^{(j)})^\top S^{(j)}\Phi)^{-1}\Phi^\top(S^{(j)})^\top S^{(j)}y
    \]
    and our final ensemble predictor is:
    \[
        \hat{\theta} = \frac{1}{m} \sum_{j=1}^m \hat{\theta}^{(j)}
    \]
    By rotating the geometry using the \example{\textbf{Singular Value Decomposition (SVD)}} of \(\Phi\), and relying on the rotational invariance of Gaussian matrices, we can write the prediction vector \(\Phi \hat \theta^{(j)}\) as \\
    \(\Phi \hat \theta^{(j)} = \begin{pmatrix} y_1 + ((S_1^{(j)})^\top S_1^{(j)})^{-1}(S_1^{(j)})^\top S_2^{(j)} y_2 \\ 0 \end{pmatrix}\)
\end{frame}

\subsection{Bias and Variance Analysis}
\begin{frame}
    Taking expectations over the random sketching matrix \(S^{(j)}\) and using the fact that \(\mathbb{E}[S_2^{(j)}] = 0\), we find that:
    \[
        \mathbb{E}_{S^{(j)}}[\Phi\hat{\theta}^{(j)}] = \Phi\hat{\theta}_{\text{OLS}}
    \]
    This implies that Gaussian Sketching provides an \alert{\textbf{unbiased}} estimator of the \alert{\textbf{true OLS prediction}}.

    However, sketching injects algorithmic variance. Utilizing the trace of the \emphasize{\textbf{Inverse Wishart distribution}}, the variance of a single sketch is:
    \[
        \mathbb{E}_{\varepsilon, S^{(j)}}\left[ \big\|\Phi\hat{\theta}^{(j)} - \mathbb{E}_{S^{(j)}}\Phi\hat{\theta}^{(j)}\big\|_2^2 \right] = \sigma^2 \frac{d(n-d)}{s-d-1}
    \]
\end{frame}

\subsection{Expected Generalization Error}
\begin{frame}
    Combining the Bias and Variance terms for our ensemble of \(m\) independent sketches, the expected fixed-design error is\footnote{See Section 3.5 for standard OLS baseline derivations.}:
    \[
        \frac{1}{n} \mathbb{E}_{\varepsilon, S} \Bigg[ \bigg\| \frac{1}{m} \sum_{j=1}^m \Phi\hat{\theta}^{(j)} - \Phi\theta_* \bigg\|_2^2 \Bigg] = \frac{\sigma^2 d}{n} + \textcolor{red}{\frac{\sigma^2 d}{n}\left(\frac{1}{m}\frac{n-d}{s-d-1}\right)}
    \]
    where:
    \begin{itemize}
        \item \(\frac{\sigma^2 d}{n}\) is the \alert{\textbf{standard OLS error}} baseline;
        \item the \textcolor{red}{red term} is the \alert{\textbf{algorithmic sketching penalty}};
        \item \(m\) is the \alert{\textbf{number of replications}} (ensemble size);
        \item \(s\) is the \alert{\textbf{reduced sketching dimension}} (\(s > d\));
        \item \(\sigma^2\) is the variance of the independent zero-mean noise \(\varepsilon\).
    \end{itemize}
\end{frame}

\begin{frame}
    What does this bound tell us about the trade-offs?
    \begin{itemize}
        \item as \example{\(m \to \infty\)}: the penalty vanishes. Averaging infinitely many sketches perfectly reconstructs the original OLS estimator.
        \item to \example{\textbf{match OLS performance}} (up to a factor of 2), we need the penalty term to equal the base error. This requires \(m \sim \frac{n}{s}\) replications.
    \end{itemize}
    \vspace{0.5cm}
    \textbf{Conclusion:} Just like bagging, there is \emphasize{\textbf{no statistical gain}} over OLS. The benefit is purely computational. We can process \(m\) tiny datasets of size \(s \times d\) in parallel instead of one massive dataset of size \(n \times d\).
\end{frame}

\makepart{Random Projections and Extensions}
\section{Random Projections}
\subsection{The High-Dimensional Setting}
\begin{frame}
    We now flip the regime: assume \(d > n\) (high-dimensional set-up). Standard OLS is ill-posed because \(\Phi^\top\Phi\) is not invertible.
    
    For each \(j = 1,\dots,m\), we use a sketching matrix \(S^{(j)} \in \mathbb{R}^{d \times s}\) (\(s \le n\)), sambled from a distribution to be determined we just assume that \(\text{rank}(S^{(j)} = j)\) a.s. We restrict our parameters by considering \(\hat{\theta}^{(j)} = S^{(j)}\hat{\eta}^{(j)}\) where \(\hat{\eta}^{(j)}\) minimizes \(\min_{\eta \in \mathbb{R}^s} \|y - \Phi S^{(j)}\eta\|_2^2\).
    The denoised response vector is \(\hat{y}^{(j)} = \Pi^{(j)} y\), where \(\Pi^{(j)} = \Pi^{(j)} = \Phi S^{(j)} \big((S^{(j)})^\top\Phi^\top\Phi S^{(j)}\big)^{-1} (S^{(j)})^\top\Phi^\top\) is an \alert{\textbf{orthogonal projection matrix}} onto an \(s\)-dimensional vector space, such that \(\text{tr}(\Delta) = s\). Let its expectation be \(\Delta = \mathbb{E}[\Pi^{(j)}]\).
\end{frame}

\subsection{Expected Excess Risk Bound}
\begin{frame}
    Decomposing the total error into bias squared and viariance we obtain a first bound of the form \\
    \(\frac{1}{n} \mathbb{E}_{\varepsilon, S} [ \| \frac{1}{m} \sum_{j=1}^m \hat{y}^{(j)} - \Phi\theta_* \|_2^2 ] \le \frac{\sigma^2 s}{n} + \frac{1}{n} \theta_*^\top \Phi^\top [I - \Delta] \Phi\theta_*.\) \\
    Now using the properties of the expectation of projection matrix to bound \(I - \Delta\) we get 
    \begin{equation}
        \frac{1}{n} \mathbb{E}_{\varepsilon, S} [ \| \frac{1}{m} \sum_{j=1}^m \hat{y}^{(j)} - \Phi\theta_* \|_2^2 ] \le  \frac{\sigma^2 s}{n} + \frac{1}{n} \theta_*^\top \Phi^\top \Big( I - \Phi \mathbb{E}_S \big[SS^\top\big] \big( \mathbb{E}_S \big[SS^\top \Phi^\top \Phi SS^\top\big] \big)^{-1} \mathbb{E}_S \big[SS^\top\big] \Phi^\top \Big) \Phi\theta_*.
    \end{equation}
    
    Finally assuming the case of Gaussian random projections, the overall expected excess risk is bounded by:
    \[
        \underbrace{\frac{\sigma^2 s}{n}}_{\text{variance upper bound}} + \underbrace{\frac{2}{n}\text{tr}(\Phi^\top\Phi)\frac{\|\theta_*\|_2^2}{s+1}}_{\text{bias upper bound}}
    \]
    where:
    \begin{itemize}
        \item \(s\) acts as the \alert{\textbf{inverse regularization parameter}} (similar to \(1/\lambda\) in Ridge Regression)
    \end{itemize}
    As \(s\) drops, variance vanishes but bias explodes. This perfectly mirrors the \example{\textbf{bias-variance trade-off}} seen in Ridge Regression (Eq. 3.6).
\end{frame}

\subsection{Exercise 10.4: Feature Subsampling}

\begin{frame}
    What if we replace the dense Gaussian matrix with a sparse one? 

    Consider a sketching matrix \(S \in \mathbb{R}^{d \times s}\), where each column is a \alert{\textbf{canonical basis vector}} \(e_i\). This is identical to randomly \example{\textbf{subsampling \(s\) features}} out of \(d\).

    Using the the bias bound from eq. (1) and using the moments of the Multinomial distribution we get an \example{\textbf{overall expected excess risk}}:
    \[
        \underbrace{\frac{\sigma^2 s}{n}}_{\text{variance}} + \underbrace{\frac{1}{n} \left( \frac{d}{s-1} \text{tr}(\Phi^\top\Phi) \|\theta_*\|_2^2 \right)}_{\text{bias upper bound}}
    \]

    \textbf{Conclusion:}
    \begin{itemize}
        \item The bias term scales as \(\frac{d}{s}\) instead of \(\frac{1}{s}\) (Gaussian case). You need \(s \propto d\) to achieve the same regularization.
        \item However, it provides better \alert{\textbf{computational performance}} since no dense \(\mathcal{O}(nd^2)\) matrix multiplications are required.
    \end{itemize}
\end{frame}


\section{Extensions: Kernel Methods}

\subsection{Infinite-Dimensional Models and RKHS}
\begin{frame}
    How do we apply random projections when the feature dimension \(d\) is infinite? We turn to \emphasize{\textbf{Kernel Methods}}
    
    Instead of explicit feature vectors \(\varphi(x)\), we access data exclusively through a \alert{\textbf{positive definite kernel function}}:
    \[
        k(x, x') = \langle \varphi(x), \varphi(x') \rangle_{\mathcal{H}}
    \]
    where \(\mathcal{H}\) is a \emphasize{\textbf{Reproducing Kernel Hilbert Space (RKHS)}}. Key properties from Chapter 7 include:
    \begin{itemize}
        \item \alert{\textbf{Reproducing Property}}: Function evaluation is a dot-product: \(\langle k(\cdot, x), f \rangle_{\mathcal{H}} = f(x)\).
        \item \alert{\textbf{Representer Theorem}}: The infinite-dimensional optimization \(\min_{f \in \mathcal{H}} \frac{1}{n} \sum \ell(y_i, f(x_i)) + \frac{\lambda}{2} \|f\|_{\mathcal{H}}^2\) has a finite-dimensional solution \(f(x) = \sum_{i=1}^n \alpha_i k(x, x_i)\).
    \end{itemize}
    This allows us to map the target signal \(y_* = \Phi\theta_*\) to \(K\alpha\), and the norm \(\|\theta_*\|_2^2\) to the RKHS norm \(y_*^\top K^{-1} y_*\).
\end{frame}

\subsection{Bochner's Theorem and the Generalization Bound}
\begin{frame}
    To apply our random projection theory to kernels, we consider a "sketch" \(\hat{\Phi} \in \mathbb{R}^{n \times s}\). What defines the underlying geometry of these projections?
    
    \emphasize{\textbf{Bochner's Theorem (Prop 7.4):}} A translation-invariant kernel \(k(x, x') = q(x - x')\) is positive definite if and only if \(q\) is the \example{\textbf{Fourier transform of a non-negative Borel measure}}.
    \[
        k(x, x') = \int_{\mathbb{R}^d} \hat{q}(\omega) e^{i\omega^\top(x-x')} d\omega
    \]
    If we mathematically sample the columns of our sketch \(\hat{\Phi}\) from a multivariate Gaussian \(\mathcal{N}(0, K)\), we perfectly map the ensemble projection bounds from Chapter 10 into the kernel space:
    \[
        \text{Expected Risk} \le \underbrace{\frac{\sigma^2 s}{n}}_{\text{variance}} + \underbrace{\frac{2}{n}\text{tr}(K)\frac{y_*^\top K^{-1} y_*}{s+1}}_{\text{bias (RKHS norm)}}
    \]
    We recover the exact same \example{\textbf{Bias-Variance trade-off}}, dictated by the sketch dimension \(s\).
\end{frame}

\subsection{Random Fourier Features (RFF)}
\begin{frame}
    \alert{\textbf{The Bottleneck:}} Generating a sketch \(\hat{\Phi} \sim \mathcal{N}(0, K)\) requires computing the full \(n \times n\) matrix \(K\) and its square root \(K^{1/2}\). This takes \example{\(\mathcal{O}(n^3)\) time}, destroying the computational benefit of random projections!
    
    \alert{\textbf{The Solution:}} We use \emphasize{\textbf{Random Fourier Features}}. Using Bochner's theorem, we rewrite the kernel as an expectation over random frequencies \(v\):
    \[
        k(x, x') = \mathbb{E}_v \big[ \varphi(x, v)\varphi(x', v) \big]
    \]
    Instead of computing \(K\), we construct \(\hat{\Phi} \in \mathbb{R}^{n \times s}\) directly. Each column evaluates a random frequency across all datapoints: \(\big(\varphi(x_1, v), \dots, \varphi(x_n, v)\big)^\top\).
    \begin{itemize}
        \item \example{\textbf{Time to build sketch:}} \(\mathcal{O}(ns)\) (no \(n \times n\) matrix computed).
        \item \example{\textbf{Time to solve OLS:}} \(\mathcal{O}(ns^2)\).
    \end{itemize}
    RFF gives us a highly efficient, strictly \(\mathcal{O}(n)\) algorithm that practically achieves the theoretical infinite-dimensional risk bounds!
\end{frame}

\subsection{The Johnson-Lindenstrauss Lemma}
\begin{frame}
    Why do random projections work so well geometrically? 
    
    The classical \alert{\textbf{Johnson-Lindenstrauss (JL) Lemma}} states that for any \(n\) points in \(\mathbb{R}^d\), a Gaussian random projection matrix \(S \in \mathbb{R}^{d \times s}\) preserves pairwise distances with probability \(> 1-\delta\) if:
    \[
        s \ge \frac{6}{\varepsilon^2} \log \frac{n^2}{\delta}
    \]
    yielding:
    \[
        (1-\varepsilon)\|\varphi_i - \varphi_j\|_2^2 \le \|s^{-1/2}S^\top \varphi_i - s^{-1/2}S^\top \varphi_j\|_2^2 \le (1+\varepsilon)\|\varphi_i - \varphi_j\|_2^2
    \]
    Crucially, the required projected dimension \(s\) depends \example{\textbf{only logarithmically on \(n\)}} and is completely independent of the original dimension \(d\).
\end{frame}


\end{document}
